微处理器实现中可能存在规范未显式描述、工具链难以一致识别、或在不同执行环境下表现不一致的指令行为，这类差异会直接影响系统软件验证、漏洞分析与第三方审计。现有工作已经证明隐藏指令审计具有研究价值，但在多 ISA 统一工程结构、真实硬件证据链、以及新架构低成本接入方面仍存在明显缺口。本文设计并实现了一种面向 RISC 风格处理器的隐藏指令分析工具，以单条指令注入为核心，将 \memcage\ 与 \ptracebackend\ 两类执行后端、信号与反汇编联合分类、以及配置驱动的新架构后端生成纳入同一框架。当前结果表明，该工具已在 Layer A 环境下稳定完成 RISC-V 与 AArch64 的扫描与日志落盘；在 AArch64 对齐子集上，\memcage\ 的 median 吞吐达到 9362.29 insn/s，高于 \ptracebackend\ 的 8192.00 insn/s；同时，MIPS64 生成链路已完成 quick smoke，说明“配置 + 生成 + 有界人工加固”的路线具备可行性。本文当前稿的边界在于 Layer B 实机验证尚未完成，因此关于实机可复现性与 Layer A/Layer B 差异归因的结论仍保留为占位，但现有实现已经形成了一套可扩展、可复现、可审计的分析框架，并为后续板卡补证提供了直接落点。
